%
% Documento: Introdução
%
\chapter{INTRODUÇÃO}\label{chap:introducao}

Sistemas distribuídos são formados por processos independentes que cooperam por
troca de mensagens. Como não há memória nem relógio físico compartilhados, a
observação de eventos e a manutenção de uma ordem comum exigem protocolos
explícitos. Relógios lógicos preservam relações de precedência, enquanto relógios
vetoriais permitem representar relações causais e identificar eventos concorrentes
\cite{lamport1978,tanenbaum2007}. Ainda assim, causalidade define apenas uma ordem
parcial: mensagens concorrentes podem ser observadas em ordens diferentes por
participantes distintos.

O trabalho propõe a implementação de uma aplicação distribuída com comunicação de
grupo, ordem total, número configurável de nós, interface por nó e captura de um
estado global consistente. Entre os temas oferecidos, escolheu-se o
\textbf{chat distribuído}. Cada participante é um processo do sistema operacional,
com estado privado e endereço próprio. A única informação compartilhada antes da
execução é o catálogo estático de endereços; durante a execução, toda coordenação
ocorre por mensagens de rede.

Para estabelecer uma sequência única das mensagens de grupo, foi escolhida a
abordagem de \textbf{sequenciador}. Um coordenador atribui números globais
monotônicos e redifunde as mensagens; relógios vetoriais permanecem anexados aos
eventos para registrar a evolução causal. Como o coordenador é um ponto crítico,
o protótipo emprega batimentos de saúde e o algoritmo do Valentão para eleição de
um substituto. Essa opção reduz a complexidade do caminho normal de entrega, ao
custo de exigir tratamento cuidadoso para a falha do líder.

O objetivo deste relatório é descrever fielmente a arquitetura encontrada no
repositório, explicar seus algoritmos e registrar os resultados verificáveis. Por
isso, além das funcionalidades demonstradas, são apresentadas as limitações ainda
existentes. Em especial, a versão analisada não implementa a captura de estado
global e não conclui o recebimento de mensagens privadas, embora a interface de
envio já exista.

\section{Organização do relatório}

O Capítulo~\ref{chap:projeto} descreve equipe, tecnologias, arquitetura, rede e
execução. O Capítulo~\ref{chap:algoritmos} detalha relógio vetorial, ordem total,
eleição e a alternativa planejada para estado global. O
Capítulo~\ref{chap:avaliacao} apresenta os cenários de validação e as limitações.
Por fim, o Capítulo~\ref{chap:conclusao} sintetiza os resultados e os requisitos
pendentes.
